\ifdefined\iscombined
	\lecturetitle{Irreducible Quadratics and Stability}
\else
\documentclass{article}
\usepackage{xparse}
\usepackage{changepage}
\usepackage{listings}
\usepackage{amsthm}
\usepackage{comment}
\usepackage{amsmath}
\usepackage{braket}
\usepackage{amsfonts}
\usepackage{sectsty}
\usepackage{hyperref}
\usepackage{fancyhdr}
\usepackage[top=1in,bottom=1in,left=1in,right=1in]{geometry}
\usepackage{float}
\usepackage{graphicx}
\usepackage{mathtools}
\usepackage{tikz}
\usetikzlibrary{arrows}

\date{
	\vspace{-.75em}
	{\large \today}
}

\title{
	\vspace*{-1.5em}
	{\Huge Lecture 4} \\
	\vspace{-1em}
	\rule{\textwidth/4}{.01em} \\
	\vspace{-.25em}
	{\Large Irreducible Quadratics and Stability}
	\vspace{-.75em}
}

\author{
	{\large Matthew Farah}
}

\hypersetup{
	colorlinks=true,
	linkcolor=blue,
	filecolor=magenta,
	urlcolor=blue,
	citecolor=blue,
	anchorcolor=blue
}

\fancyhead{}
\fancyhead[L]{\today}
\fancyhead[R]{Lecture 4 - Irreducible Quadratics and Stability}

\setlength{\parindent}{0cm}

\newcounter{example}
\renewcommand{\theexample}{\arabic{example}}

\newenvironment{example}[1][]{
	\refstepcounter{example}
	\begin{adjustwidth}{1.5em}{}
	\noindent\textbf{\ifx&#1&Example \theexample:\else Example \theexample \ (#1):\fi}
	\ignorespaces
}{
	\vspace{.5em}
	\end{adjustwidth}
}

\newenvironment{solution}{
	\vspace{.5em}
	\par
	\begin{adjustwidth}{1.5em}{}
	\textbf{{Solution:}}
	\par
	\vspace{.5em}
}{
	\vspace{1em}
	\end{adjustwidth}
}

\newcounter{subexample}
\renewcommand{\thesubexample}{\alph{subexample}}

\newenvironment{subexample}{
	\setcounter{step}{0}
	\refstepcounter{subexample}
	\vspace{1em}\par\noindent
	\begin{adjustwidth}{1.5em}{}
	\noindent\textbf{(\thesubexample)}
	\ignorespaces
}{
	\par
	\vspace{1em}
	\end{adjustwidth}
}

\renewenvironment{proof}{
	\vspace{.5em}\par\noindent
	\begin{adjustwidth}{1.5em}{}
	\emph{Proof:}\par\vspace{1em}
	\setlength{\parindent}{0cm}
}{
	\end{adjustwidth}
}

\newtheorem{theorem}{Theorem}

\renewenvironment{theorem}[1][]{
	\refstepcounter{theorem}
	\vspace{1em}\par\noindent
	\begin{adjustwidth}{1.5em}{}
	\noindent\textbf{\ifx&#1&Theorem \thetheorem:\else Theorem \thetheorem \ (#1):\fi}
	\ignorespaces
}{
	\vspace{.5em}
	\end{adjustwidth}
}

\newtheorem{lemma}{Lemma}

\renewenvironment{lemma}[1][]{
	\refstepcounter{lemma}
	\vspace{1em}\par\noindent
	\begin{adjustwidth}{1.5em}{}
	\noindent\textbf{\ifx&#1&Lemma \thelemma:\else Lemma \thelemma \ (#1):\fi}
	\ignorespaces
}{
	\vspace{.5em}
	\end{adjustwidth}
}

\makeatother
\makeatletter

\newcounter{step}
\renewcommand{\thestep}{\Roman{step}}

\newenvironment{step}{
	\refstepcounter{step}
	\vspace{1em}\par\noindent
	\ignorespaces\noindent\textbf{(\thestep)}
	\ignorespaces
}{
	\par
}

\sectionfont{\fontsize{12}{12}\bfseries}
\subsectionfont{\fontsize{10}{12}\normalfont}
\subsubsectionfont{\fontsize{10}{12}\normalfont}

\begin{document}

\maketitle
\pagestyle{fancy}

\fi

\begin{itemize}
	\item Note: the midterm will have a physical modelling question for mechanical and electrical systems.
	\item A pole is a value of $s$ which makes the denominator of a transfer function zero.
	\begin{itemize}
		\item A pole is said to be `unstable' if its real component is greater than zero, and `semi-stable' if its real component equals zero.
	\end{itemize}
	\item A quadratic is said to be `irreducible' if its roots are complex numbers with non-zero imaginary components.
\end{itemize}

\begin{example}
	Compute $y(t)$ where $Y(s) = \frac{3}{s(s^2 + 2s + 5)}$.
\end{example}

\begin{solution}
	Note that $Y(s)$ has an irreducible quadratic in the denominator. As earlier, to compute $y(t)$, we must decompose $Y(s)$.

	\begin{step}
		Decompose $Y(s)$.
	\end{step}

	\begin{align*}
		\frac{3}{s(s^2 + 2s + 5)} &= \frac{A}{s} + \frac{Bs + C}{s^2 + 2s + 5} \\
	\end{align*}

	We therefore have that:

	\begin{align*}
		3 &= A(s^2 + 2s + 5) + (Bs + C)(s) \\
		3 &= As^2 + 2As + 5A + Bs^2 + Cs \\
	\end{align*}

	Thus, by equating coefficients, we have:

	\begin{align*}
		5A &= 3 \Rightarrow A = \tfrac{3}{5} \\
		As^2 + Bs^2 &= 0 \Rightarrow Bs^2 = -\tfrac{3}{5}s^2 \Rightarrow B = -\tfrac{3}{5} \\
		2As + Cs &= 0 \Rightarrow Cs = -\tfrac{6}{5}s \Rightarrow C = -\tfrac{6}{5} \\
	\end{align*}

	Therefore:

	\begin{align*}
		Y(s) &= \frac{3}{5s} + \frac{-3s - 6}{5(s^2 + 2s + 5)} \\
	\end{align*}

	\begin{step}
		Compute $y(t)$.
	\end{step}

	We know that:

	\begin{align*}
		y(t) &= \mathcal{L}^{-1}\{Y(s)\} \\
			&= \mathcal{L}^{-1} \left\{ \frac{3}{5s} - \frac{3s + 6}{5(s^2 + 2s + 5)} \right\} \\
			&= \frac{3}{5} \mathcal{L}^{-1} \left\{ \frac{1}{s} \right\} - \frac{3}{5} \mathcal{L}^{-1} \left\{ \frac{s + 2}{s^2 + 2s + 5} \right\} \\
	\end{align*}

	$\mathcal{L}^{-1} \left\{ \frac{1}{s} \right\}$ is easy enough to find, but what about $\mathcal{L}^{-1} \left\{ \frac{s + 2}{s^2 + 2s + 5} \right\}$? Well, through some clever algebra, observe that:

	\begin{align*}
		\frac{s + 2}{s^2 + 2s + 5} &= \frac{s + 2}{s^2 + 2s + 1 + 4} \\
			&= \frac{s + 2}{(s + 1)^2 + 4} \\
			&= \frac{s + 1 + 1}{(s + 1)^2 + 2^2} \\
			&= \frac{s + 1}{(s + 1)^2 + 2^2} + \frac{1}{(s + 1)^2 + 2^2} \\
			&= \frac{s + 1}{(s + 1)^2 + 2^2} + \frac{1}{2} \left( \frac{2}{(s + 1)^2 + 2^2} \right) \\
	\end{align*}

	Now, note that $\mathcal{L}\{e^{-at}\cos(\omega t)\} = \frac{s + a}{(s + a)^2 + \omega^2}$ and $\mathcal{L}\{e^{-at}\sin(\omega t)\} = \frac{\omega}{(s + a)^2 + \omega^2}$. This is an awful lot like what we found earlier! Thus:

	\begin{align*}
		y(t) &= \frac{3}{5} \mathcal{L}^{-1} \left\{ \frac{1}{s} \right\} - \frac{3}{5} \mathcal{L}^{-1} \left\{ \frac{s + 1}{(s + 1)^2 + 2^2} \right\} - \frac{3}{10} \mathcal{L}^{-1} \left\{ \frac{2}{(s + 1)^2 + 2^2} \right\} \\
			&= \frac{3}{5} u(t) - \frac{3}{5} e^{-t} \cos(2t) - \frac{3}{10} e^{-t} \sin(2t)
	\end{align*}
\end{solution}

\ifdefined\iscombined
	\newpage
\else
	\end{document}
\fi
